\documentclass[conference]{IEEEtran}

\usepackage[utf8]{inputenc}
\usepackage[T1]{fontenc}
\usepackage{cite}
\usepackage{amsmath,amssymb}
\usepackage{graphicx}
\usepackage{booktabs}
\usepackage{url}
\usepackage{hyperref}
\usepackage{xcolor}
\usepackage{listings}
\usepackage{textcomp}

\lstset{
  basicstyle=\footnotesize\ttfamily,
  breaklines=true,
  frame=single,
  numbers=left,
  numberstyle=\tiny,
  keywordstyle=\color{blue},
  commentstyle=\color{gray},
  columns=fullflexible
}

\begin{document}

\title{2-Way Set-Associative Cache with LRU Replacement:\\Phase~I Prototype and Comparative Evaluation}

\author{
\IEEEauthorblockN{Krishna Pramod Dubale, Apoorva Choudhary, Prachee Dhar, Kushal Sarkar, Vani Rudra}
\IEEEauthorblockA{Modern Computer Architecture (MCA) --- Project~7\\
Rishihood University}
}

\maketitle

% ============================================================
\begin{abstract}
We present the Phase~I prototype of a synthesizable 2-Way Set-Associative
Cache with Least-Recently-Used (LRU) replacement policy, implemented in
Verilog-2001 and verified using Icarus Verilog.  The cache provides 128
bytes of data storage organized as 4~sets with 2~ways per set and 16-byte
cache blocks.  A direct-mapped cache baseline of identical capacity serves
as a comparison target.  Both designs share a common external interface with
synchronous active-high reset, a single outstanding request handshake
protocol, and a combinational lower-memory refill port.  Automated
simulation across five test suites totaling 618~assertions confirms
functional correctness.  On a 5-access conflict trace targeting a single
set, the 2-way design achieves a 40\% hit rate compared to the
direct-mapped baseline's 0\% hit rate, demonstrating the effectiveness of
LRU set associativity in mitigating ping-pong thrashing.  A spatial
locality trace confirms both architectures achieve identical 66.7\% hit
rates through multi-word block reuse, validating that the performance
advantage is attributable to associativity rather than capacity.
\end{abstract}

\begin{IEEEkeywords}
cache memory, set associativity, LRU replacement, direct-mapped cache,
conflict miss, Verilog, hardware verification
\end{IEEEkeywords}

% ============================================================
\section{Introduction}

Cache memories occupy a critical position in modern computer architecture,
bridging the performance gap between fast processors and slower main
memory~\cite{patterson2017}.  Direct-mapped caches offer simplicity and
fast access times but suffer from conflict misses when multiple frequently
accessed addresses map to the same set~\cite{hill1988}.  Set-associative
caches mitigate this by providing multiple storage locations (ways) per set,
reducing conflict misses at the cost of additional comparator hardware and a
replacement policy~\cite{hennessy2019}.

This project implements and compares a 2-way set-associative cache with LRU
replacement against a direct-mapped baseline, both with 128~bytes of data
storage.  The Phase~I deliverable provides working read-path prototypes with
automated verification evidence.

% ============================================================
\section{Literature Review}

Smith~\cite{smith1982} provided one of the earliest comprehensive surveys
of cache memory design, establishing the taxonomy of compulsory, capacity,
and conflict misses that remains foundational to cache analysis.  His work
demonstrated that increasing associativity reduces conflict misses, with
diminishing returns beyond 4-way or 8-way configurations.

Hill~\cite{hill1988} argued that direct-mapped caches could achieve
performance competitive with set-associative designs under typical
workloads, primarily because their shorter access times compensate for
higher miss rates.  However, Hill's analysis acknowledged that
pathological access patterns---where multiple hot addresses map to the
same set---expose the direct-mapped cache to severe thrashing.

Jouppi~\cite{jouppi1990} proposed augmenting direct-mapped caches with
small fully-associative victim caches to capture conflict misses.  This
work implicitly demonstrated that conflict misses represent a significant
performance concern, motivating the set-associative designs that the
victim cache seeks to approximate.

Przybylski et al.~\cite{przybylski1988} conducted systematic performance
tradeoff analyses across cache sizes, associativities, and block sizes,
concluding that 2-way set associativity captures the majority of the
conflict-miss reduction benefit for a given capacity, making it a
cost-effective design point.

Al-Zoubi et al.~\cite{al-zoubi2004} evaluated cache replacement policies
including LRU, FIFO, and random replacement using the SPEC~CPU2000 benchmark
suite, confirming that LRU consistently outperforms other policies for
2-way and 4-way configurations.  Their empirical data supports the
selection of LRU as the replacement policy for our 2-way design.

Patterson and Hennessy~\cite{patterson2017} formalize the Average Memory
Access Time (AMAT) model as:
\begin{equation}
\text{AMAT} = t_{\text{hit}} + r_{\text{miss}} \times p_{\text{miss}}
\label{eq:amat}
\end{equation}
where $t_{\text{hit}}$ is the hit time, $r_{\text{miss}}$ is the miss
penalty, and $p_{\text{miss}}$ is the miss rate.  Reducing the miss rate
through associativity directly improves AMAT, provided the hit time
increase from additional comparators is modest.

Our work follows these established principles to construct a minimal but
complete 2-way LRU cache and direct-mapped baseline, verified through
automated simulation using Icarus Verilog~\cite{icarus-verilog} with
Verilog-2001~\cite{ieee-verilog2001} RTL and SystemVerilog testbenches.

% ============================================================
\section{Proposed Solution}

\subsection{Design Overview}

Both cache designs share identical parameters: 128~bytes of data capacity
organized as 16-byte blocks (four 32-bit words per block), 32-bit byte
addressing, and an identical external interface for direct comparison.

\begin{itemize}
  \item \textbf{2-Way Set-Associative} (\texttt{cache\_2way}): 4~sets $\times$ 2~ways.  Tag: \texttt{addr[31:6]} (26~bits).  Set index: \texttt{addr[5:4]} (2~bits).  Word select: \texttt{addr[3:2]}.
  \item \textbf{Direct-Mapped Baseline} (\texttt{cache\_direct}): 8~sets $\times$ 1~way.  Tag: \texttt{addr[31:7]} (25~bits).  Set index: \texttt{addr[6:4]} (3~bits).  Word select: \texttt{addr[3:2]}.
\end{itemize}

\subsection{Architecture}

The 2-way cache consists of two modules.  The storage submodule
(\texttt{cache\_2way\_store}) encapsulates the dual-way tag and data
arrays, parallel tag comparators, valid bit registers, per-set 1-bit
LRU registers, a deterministic victim selection unit, and a 4-word read
multiplexer.

The top-level controller (\texttt{cache\_2way}) implements a 4-state
finite state machine (IDLE, LOOKUP, REFILL, RESP) that sequences each
accepted request through tag comparison, optional memory refill, and
response delivery.

\begin{figure}[ht]
\centering
\small
\begin{verbatim}
  31            6 5  4 3  2 1  0
 +---------------+----+----+----+
 |     TAG       |SET |WORD|BYTE|
 |   [31:6]      |[5:4|[3:2|[1:0|
 +---------------+----+----+----+
      26 bits     2b   2b   2b
                (4 sets)
\end{verbatim}
\caption{Address field decomposition for the 2-way set-associative cache (128~bytes, 4~sets $\times$ 2~ways, 16-byte blocks).}
\label{fig:addr-fields}
\end{figure}

\subsection{Block Diagram}

Fig.~\ref{fig:block} illustrates the top-level datapath.  The controller
latches \texttt{req\_addr} into \texttt{saved\_addr} upon handshake
acceptance.  The storage submodule performs combinational parallel tag
comparison on both ways.  On a hit, \texttt{touch\_en} updates the LRU
recency bit; on a miss, the memory refill port fetches a 128-bit block
and \texttt{fill\_en} writes it into the victim way selected by the LRU
or invalid-slot priority logic.

\begin{figure}[ht]
\centering
\small
\begin{verbatim}
    [Processor Master Interface]
       | req_valid, req_addr
       v
     +-------------------------+
     |     Controller FSM      | -> resp_valid
     | (IDLE/LOOKUP/FILL/RESP) | -> resp_hit
     +------------+------------+ -> resp_data
                  |
            saved_addr
                  |
       +----------+----------+
       |                     |
       v                     v
   +--------+           +-----------+
   | Lower  |           | cache_2way|
   | Memory | 128-bit   | _store    |
   | Refill |---------->| Way0/Way1 |
   | Port   |           | LRU & MUX |
   +--------+           +-----------+
\end{verbatim}
\caption{Simplified block diagram of \texttt{cache\_2way}.  The controller FSM orchestrates the handshake, lookup, refill, and response phases.}
\label{fig:block}
\end{figure}

\subsection{Controller FSM}

The finite state machine operates as follows:
\begin{enumerate}
  \item \textbf{IDLE}: \texttt{req\_ready}~=~1.  On \texttt{req\_valid $\land$ req\_ready}, latch address, transition to LOOKUP.
  \item \textbf{LOOKUP}: Evaluate tag match.  On hit, assert \texttt{touch\_en} for one cycle, register response, transition to RESP.  On miss, transition to REFILL.
  \item \textbf{REFILL}: Assert \texttt{fill\_en} for one cycle, writing \texttt{mem\_block\_data} into the victim way.  Register response with \texttt{resp\_hit}~=~0, transition to RESP.
  \item \textbf{RESP}: Deliver one-cycle \texttt{resp\_valid} pulse.  Return to IDLE.
\end{enumerate}

The controller enforces a mutual exclusion invariant: \texttt{touch\_en} and \texttt{fill\_en} are never asserted simultaneously.

\subsection{LRU Replacement and Worked Example}

Each set maintains a 1-bit LRU register.  When \texttt{lru\_bit}~=~0, Way~0 is the least recently used; when \texttt{lru\_bit}~=~1, Way~1 is the LRU victim.  Invalid slots receive absolute priority during allocation.

Consider the 5-access trace starting from a cold cache:
\begin{center}
\texttt{0x00}, \texttt{0x80}, \texttt{0x00}, \texttt{0x100}, \texttt{0x00}
\end{center}

All three unique addresses map to Set~0 in both designs.  In the 2-way cache:

\begin{enumerate}
  \item \textbf{Request~1} (\texttt{0x00}): Miss.  Fill Way~0.  \texttt{lru\_bit}~$\leftarrow$~1.
  \item \textbf{Request~2} (\texttt{0x80}): Miss.  Fill Way~1.  \texttt{lru\_bit}~$\leftarrow$~0.
  \item \textbf{Request~3} (\texttt{0x00}): \textbf{Hit} on Way~0.  \texttt{touch\_en} flips \texttt{lru\_bit} to~1, protecting \texttt{0x00} from eviction.
  \item \textbf{Request~4} (\texttt{0x100}): Miss.  LRU selects Way~1 (\texttt{0x80}) as victim.  \texttt{0x00} survives in Way~0.
  \item \textbf{Request~5} (\texttt{0x00}): \textbf{Hit} on Way~0.  Confirmed at $t=196$\,ns in waveform evidence.
\end{enumerate}

Result: 2/5 hits (40\%).  The direct-mapped baseline achieves 0/5 hits (0\%) because each access to Set~0 evicts the sole resident block.

\subsection{Verification Methodology}

All RTL is written in synthesizable Verilog-2001 and compiled using Icarus Verilog with the \texttt{-g2012} flag for SystemVerilog testbenches~\cite{icarus-verilog}.  A deterministic testbench memory model supplies block data where each 32-bit word equals its own aligned byte address, enabling data integrity verification independent of hit/miss classification.

Five automated test suites verify the design:
\begin{itemize}
  \item Storage unit tests (51~assertions): Cold state, lookup invariance, empty-way priority, word multiplexing, LRU touch, eviction, set independence, and reset.
  \item 2-way integration tests (141~assertions): Protocol, 5-request demo trace, spatial locality, address mutation immunity, idle stability, and reset.
  \item Direct-mapped tests (30~assertions): Protocol, conflict trace, spatial hits, mutation immunity, and reset.
  \item Comparative trace harness (52~assertions): Dual-cache concurrent execution under identical conflict and locality traces.
  \item System regressions (344~assertions): Full-population reset, post-reset cold access, back-to-back streaming, and cross-set LRU independence.
\end{itemize}

All 618~assertions pass with zero failures.

% ============================================================
\section{Initial Comparison Results}

Table~\ref{tab:comparison} summarizes the verified comparison results from automated simulation.

\begin{table}[ht]
\centering
\caption{Verified Simulation Results: 2-Way vs.\ Direct-Mapped Cache}
\label{tab:comparison}
\begin{tabular}{@{}lccccc@{}}
\toprule
\textbf{Trace} & \textbf{Accesses} & \multicolumn{2}{c}{\textbf{2-Way}} & \multicolumn{2}{c}{\textbf{Direct}} \\
\cmidrule(lr){3-4} \cmidrule(lr){5-6}
 & & Hits & Rate & Hits & Rate \\
\midrule
A (Conflict)  & 5 & 2 & 40\%   & 0 & 0\%  \\
B (Locality)  & 6 & 4 & 66.7\% & 4 & 66.7\% \\
\bottomrule
\end{tabular}
\end{table}

Trace~A demonstrates the conflict-miss advantage of set associativity: the 2-way cache retains a working set of two blocks in a single set, whereas the direct-mapped cache suffers complete thrashing.  Trace~B confirms that the spatial locality benefit of multi-word blocks is shared equally by both architectures, isolating the performance gain to associativity rather than capacity.

% ============================================================
\section{Individual Contributions}

\begin{enumerate}
  \item \textbf{Kushal Sarkar} (Scaffold and Verification, Handoffs~1 and~4):
    Initialized the repository, authored the project brief
    (\texttt{PROJECT\_BRIEF.md}), interface contracts
    (\texttt{docs/interface.md}), benchmark trace analysis
    (\texttt{docs/demo-trace.md}), and project architecture contracts.
    In Round~2, implemented system regression
    testbenches (\texttt{tb/tb\_system\_regression.sv}, 344~assertions),
    the master verification workflow (\texttt{scripts/run\_all.sh}),
    and the reproducible demo waveform evidence
    (\texttt{docs/evidence/cache\_2way\_demo.vcd}).
    Commits: \texttt{3043e05}, \texttt{92dfa5b}, \texttt{3eb612c},
    \texttt{d6cdad6}, \texttt{005cb92}, \texttt{fdb01a1}.

  \item \textbf{Krishna Pramod Dubale} (Storage and LRU, Handoff~2):
    Implemented the synthesizable \texttt{cache\_2way\_store} submodule
    (\texttt{rtl/cache\_2way\_store.v}) with dual-way tag/data arrays,
    parallel tag comparators, per-set 1-bit LRU registers, and
    deterministic victim selection.  Authored 51-assertion unit testbench
    (\texttt{tb/tb\_cache\_2way\_store.sv}) and LRU walkthrough
    documentation (\texttt{docs/lru\_explanation.md}).
    Commits: \texttt{d76542e}, \texttt{0c1a380}, \texttt{fb36862}.

  \item \textbf{Apoorva Choudhary} (Controller, Handoff~3):
    Implemented the top-level \texttt{cache\_2way} controller
    (\texttt{rtl/cache\_2way.v}) with 4-state FSM, request latching,
    LRU touch and refill datapath integration, and true miss
    classification preservation.  Authored the reusable deterministic
    memory model (\texttt{tb/mem\_block\_model.v}), 141-assertion
    integration testbench (\texttt{tb/tb\_cache\_2way.sv}), and
    architectural documentation (\texttt{docs/architecture.md}).
    Commits: \texttt{7c777f7}, \texttt{ba299cf}, \texttt{9a976c8}.

  \item \textbf{Prachee Dhar} (Comparison, Handoff~4):
    Implemented the direct-mapped cache baseline
    (\texttt{rtl/cache\_direct.v}), 30-assertion baseline testbench
    (\texttt{tb/tb\_cache\_direct.sv}), and the 52-assertion comparative
    dual-cache trace harness (\texttt{tb/tb\_cache\_compare.sv}).
    Authored verified comparison evidence
    (\path{docs/evidence/cache_comparison_evidence.md}).
    Commits: \texttt{fd3d8d7}, \texttt{7188fc6}, \texttt{cc045d7}.

  \item \textbf{Vani Rudra} (Report and Demo Package, Handoff~5):
    Prepared this IEEE Phase~I progress report (\texttt{report/main.tex},
    \texttt{report/phase1-report.pdf}), demonstration video
    materials (\texttt{demo/recording-script.md}, \texttt{demo/presentation.pdf}),
    Phase~I submission audit checklist (\texttt{docs/phase1-submission.md}), and
    final repository navigation.
    Commits: \texttt{693444d}, \texttt{da0765a}, \texttt{c0e872a},
    \texttt{f097865}, \texttt{1486e0d}, \texttt{e13598c}.
\end{enumerate}

% ============================================================
\section*{Acknowledgments}
The authors acknowledge the Modern Computer Architecture course
instructors at Rishihood University for defining the project
specifications.

% ============================================================
\bibliographystyle{IEEEtran}
\bibliography{refs}

\end{document}
